\documentclass[11pt]{article}
\usepackage[utf8]{inputenc}
\usepackage{amsmath,amssymb}
\usepackage{graphicx}
\usepackage{booktabs}
\usepackage{hyperref}
\usepackage[margin=1in]{geometry}
\usepackage{natbib}
\usepackage{multirow}

\title{DEBBIES: A Demographic Dataset for Cross-Taxonomic Comparison of Ectotherm Life History Strategies Using Dynamic Energy Budget Theory}

\author{%
  DEBBIES Consortium\\
  \texttt{correspondence@debbies-dataset.org}
}

\date{}

\begin{document}
\maketitle

\begin{abstract}
Cross-taxonomic comparison of life history strategies in ectotherms is hampered by the lack of standardized trait datasets suitable for demographic modelling. Here we present DEBBIES (Dynamic Energy Budget-Based Integrated Ectotherm Strategies), a dataset of eight standardized life history traits for 185 ectotherm species parameterized under dynamic energy budget (DEB) theory. Traits include length at birth, length at puberty, maximum length, maximum reproduction rate, energy allocation fraction, von Bertalanffy growth rate, and juvenile and adult mortality rates. These traits enable parameterization of DEB integral projection models (DEB-IPMs) for computing population growth rate, demographic resilience, and other demographic quantities. Technical validation demonstrates satisfactory agreement between model-predicted and observed longevity, generation time, and age at maturity across all 185 species. Compared to existing datasets, DEBBIES offers three key advantages: (i) standardized DEB parameterization enabling direct cross-taxonomic comparisons, (ii) inclusion of data-deficient species often excluded from traditional compilations, and (iii) support for population forecasting under novel environmental conditions via the mechanistic DEB-IPM framework. The dataset is released as an open-access resource linked to IUCN conservation status and biogeographical characteristics.
\end{abstract}

\section{Introduction}

Understanding how life history strategies vary across species is a central goal of ecology and evolutionary biology \citep{stearns1992,roff2002}. Life history traits---such as size at maturity, reproductive output, growth rate, and mortality schedules---jointly determine population dynamics and demographic responses to environmental change \citep{salguero2016}. For ectotherms, which comprise the vast majority of animal biodiversity, life history variation is particularly consequential because metabolic rate, growth, and reproduction are directly coupled to ambient temperature \citep{angilletta2009}.

Despite decades of research, cross-taxonomic comparisons of ectotherm life history strategies remain challenging. Existing trait compilations such as AmphiBIO \citep{oliveira2017}, AnAge \citep{tacutu2018}, and various taxon-specific databases use heterogeneous trait definitions, measurement methods, and environmental contexts, making direct comparison problematic. Furthermore, many ectotherm species are data-deficient, with insufficient demographic data for traditional life table analyses \citep{conde2019}.

Dynamic Energy Budget (DEB) theory provides a mechanistic framework for describing how organisms acquire and allocate energy throughout their life cycle \citep{kooijman2010}. DEB parameters have been estimated for thousands of species through the Add-my-Pet project \citep{marques2018}, providing a standardized physiological parameterization that transcends taxonomic boundaries. When coupled with integral projection models (IPMs), DEB parameters can generate complete demographic predictions---population growth rate, generation time, demographic resilience---from a small set of standardized traits \citep{smallegange2017}.

Here we present DEBBIES, a dataset of eight DEB-derived life history traits for 185 ectotherm species, designed to enable cross-taxonomic demographic modelling using DEB-IPMs. We validate the dataset by comparing model-predicted demographic quantities against observed values and demonstrate its advantages over existing compilations.

\section{Related Work}

\subsection{Life History Trait Databases}

Several databases compile life history traits for broad taxonomic groups. AmphiBIO \citep{oliveira2017} provides traits for amphibians, and AnAge \citep{tacutu2018} focuses on ageing and longevity across vertebrates. The Pantheria database \citep{jones2009} covers mammalian traits, while FishBase \citep{froese2010} and various reptile databases provide taxon-specific information. However, these databases use disparate trait definitions and measurement protocols, precluding standardized cross-taxonomic analyses.

\subsection{Dynamic Energy Budget Theory}

DEB theory \citep{kooijman2010} describes the rates at which organisms assimilate and utilize energy as a function of body size and environmental conditions. The Add-my-Pet collection \citep{marques2018} maintains DEB parameter estimates for over 3,000 species, estimated through a standardized coercion procedure. DEB-IPMs extend standard IPMs \citep{easterling2000} by using DEB parameters to mechanistically define growth, reproduction, and survival kernels, enabling population projections under novel conditions \citep{smallegange2017}.

\subsection{Demographic Modelling for Conservation}

Demographic models are increasingly used to assess extinction risk and prioritize conservation efforts \citep{salguero2016}. However, their application has been limited to data-rich species with sufficient demographic monitoring. By deriving demographic quantities from physiological parameters, DEB-IPMs can extend demographic analyses to data-deficient species, potentially informing conservation assessments for a much broader range of taxa \citep{nisbet2000}.

\section{Methods}

\subsection{Species Selection and Trait Compilation}

We compiled life history traits for 185 ectotherm species from the Add-my-Pet DEB parameter database and supplementary literature sources. Species were selected to span major ectotherm taxonomic groups (amphibians, reptiles, fish, invertebrates) and represent a range of body sizes, habitats, and conservation statuses.

\subsection{Trait Definitions}

Eight standardized traits were extracted or derived for each species (Table~\ref{tab:traits}):

\begin{table}[h]
\centering
\caption{DEBBIES trait definitions.}
\label{tab:traits}
\begin{tabular}{lll}
\toprule
Trait & Description & Units \\
\midrule
Length at birth & Size at hatching/birth & cm \\
Length at puberty & Size at reproductive maturity & cm \\
Maximum length & Asymptotic or max observed size & cm \\
Max reproduction rate & Peak reproductive output & offspring/time \\
Energy allocation fraction & Fraction to soma vs.\ reproduction & dimensionless \\
Von Bertalanffy growth rate & Somatic growth rate parameter & 1/time \\
Mortality rate (juvenile) & Pre-maturity mortality & 1/time \\
Mortality rate (adult) & Post-maturity mortality & 1/time \\
\bottomrule
\end{tabular}
\end{table}

All traits were standardized to consistent units and reference temperature (20$^\circ$C) following DEB conventions, enabling direct cross-species comparison.

\subsection{DEB-IPM Parameterization}

For each species, the eight traits were used to parameterize a DEB-IPM following the framework of \citet{smallegange2017}. The IPM defines a continuous size-structured population model where the growth kernel is derived from the von Bertalanffy growth rate and maximum length, the reproduction kernel from the reproduction rate and energy allocation fraction, and the survival kernel from the juvenile and adult mortality rates. The dominant eigenvalue of the resulting projection kernel yields the asymptotic population growth rate ($\lambda$), and additional demographic quantities (generation time, demographic resilience) are computed from the eigenvalue spectrum.

\subsection{Technical Validation}

Model-predicted longevity, generation time, and age at maturity were compared against observed values compiled from AnAge, FishBase, and primary literature. Agreement was assessed using correlation analysis across all 185 species.

\section{Results}

\subsection{Dataset Overview}

The DEBBIES dataset comprises eight standardized life history traits for 185 ectotherm species (Table~\ref{tab:overview}). Species span multiple taxonomic classes and represent a broad range of body sizes (birth lengths from $<$0.1\,cm to $>$10\,cm), growth rates, and reproductive strategies.

\begin{table}[h]
\centering
\caption{DEBBIES dataset summary.}
\label{tab:overview}
\begin{tabular}{lr}
\toprule
Property & Value \\
\midrule
Number of species & 185 \\
Taxonomic scope & Ectotherms \\
Number of traits & 8 \\
Model framework & DEB-IPM \\
\bottomrule
\end{tabular}
\end{table}

\subsection{Technical Validation}

DEB-IPM predictions showed satisfactory agreement with observed demographic quantities across all 185 species (Table~\ref{tab:validation}). Predicted longevity, generation time, and age at maturity each showed positive correlations with observed values, supporting the validity of the DEB parameterization for demographic inference. Agreement was maintained across taxonomic groups, including data-deficient species for which direct demographic observations are scarce.

\begin{table}[h]
\centering
\caption{Validation: agreement between DEB-IPM predictions and observed values.}
\label{tab:validation}
\begin{tabular}{ll}
\toprule
Demographic quantity & Agreement \\
\midrule
Longevity & Satisfactory (all 185 species) \\
Generation time & Satisfactory (cross-taxonomic) \\
Age at maturity & Satisfactory (incl.\ data-deficient spp.) \\
\bottomrule
\end{tabular}
\end{table}

\subsection{Comparison with Existing Datasets}

DEBBIES offers three key advantages over existing life history compilations (Table~\ref{tab:comparison}). First, standardized DEB parameterization enables direct cross-taxonomic comparisons without the need to harmonize heterogeneous trait definitions. Second, DEBBIES includes data-deficient species that are typically excluded from traditional compilations due to insufficient demographic monitoring data. Third, the mechanistic DEB-IPM framework supports population forecasting under novel environmental conditions (e.g., temperature change), a capability not provided by purely empirical trait databases.

\begin{table}[h]
\centering
\caption{Comparison of DEBBIES with existing datasets.}
\label{tab:comparison}
\begin{tabular}{lcc}
\toprule
Feature & DEBBIES & Existing datasets \\
\midrule
Cross-taxonomic comparison & Standardized & Difficult \\
Data-deficient species & Included & Often excluded \\
Novel-condition forecasts & DEB-IPM & Limited \\
\bottomrule
\end{tabular}
\end{table}

\section{Discussion}

The DEBBIES dataset addresses a persistent gap in comparative life history research by providing standardized, DEB-derived traits that enable rigorous cross-taxonomic demographic analyses for ectotherms. By grounding trait definitions in DEB theory, we avoid the heterogeneity that has historically complicated comparisons across taxon-specific databases.

The mechanistic basis of DEB-IPMs is a particular strength for conservation applications. Unlike purely statistical demographic models that require extensive longitudinal monitoring data, DEB-IPMs can generate demographic predictions from physiological parameters that are estimable from relatively limited biological information \citep{kooijman2010,nisbet2000}. This extends demographic analysis to data-deficient species---a critical need given that the majority of ectotherm species lack sufficient monitoring for traditional population viability analyses.

The ability to forecast population dynamics under novel conditions is especially relevant in the context of climate change. Because DEB parameters explicitly incorporate temperature dependence of metabolic processes, DEB-IPMs can project how warming temperatures will alter growth, reproduction, and survival schedules \citep{kearney2009}. This capability could inform proactive conservation strategies for thermally sensitive ectotherms.

Limitations of the current dataset include the reliance on DEB parameter estimates that themselves carry uncertainty, particularly for data-deficient species. The validation against observed demographic quantities, while satisfactory, was limited by the availability and quality of observational data for comparison. Future versions of DEBBIES will expand taxonomic coverage, incorporate parameter uncertainty estimates, and provide pre-computed demographic outputs for standard environmental scenarios.

\section{Conclusion}

We have presented DEBBIES, a standardized dataset of eight DEB-derived life history traits for 185 ectotherm species enabling cross-taxonomic demographic modelling. Technical validation confirms satisfactory agreement between DEB-IPM predictions and observed demographic quantities. By combining standardized parameterization, inclusion of data-deficient species, and mechanistic population forecasting, DEBBIES provides a resource for comparative life history research and conservation applications. The dataset is released as an open-access resource to support the ecological and conservation communities.

\bibliographystyle{plainnat}
\bibliography{references}

\end{document}
